% 分类目录：dl
\documentclass[12pt, a4paper]{article}
\usepackage[margin=2.5cm]{geometry}
\usepackage{ctex}
\usepackage{amsmath, amssymb}
\usepackage{listings}
\usepackage{xcolor}
\usepackage{tabularx}
\usepackage{hyperref}
\usepackage{tikz}
\usetikzlibrary{shapes, arrows.meta, positioning, calc}

\lstset{
    language=Python,
    basicstyle=\ttfamily\small,
    keywordstyle=\color{blue},
    commentstyle=\color{green!50!black},
    stringstyle=\color{red},
    showstringspaces=false,
    breaklines=true,
    numbers=left,
    numberstyle=\tiny\color{gray},
    frame=single
}

\title{知识点：GPT (Generative Pre-trained Transformer)}
\author{}
\date{}

\begin{document}
\maketitle

\section{核心定义}
\textbf{中文定义}：GPT（生成式预训练 Transformer）是由 OpenAI 提出的一系列自回归语言模型。其核心架构基于 Transformer 的 \textbf{Decoder-only} 结构，通过在大规模无标注文本上进行“下一个词预测”（Next Token Prediction）的生成式预训练，学习通用的语言表征能力。

\textbf{英文定义}：GPT (Generative Pre-trained Transformer) is a series of auto-regressive language models developed by OpenAI. Based on the Decoder-only Transformer architecture, it is pre-trained on massive unlabeled text via the objective of "next token prediction" to learn general-purpose language representations.

\textbf{核心哲学}：
GPT 的成功证明了“预训练 + 微调/零样本提示”范式的强大。它通过增大参数规模（Scaling Laws）和数据量，实现了从专用模型到通用人工智能（AGI）的跨越。

\section{架构：Decoder-only Transformer}

\subsection{为什么只用 Decoder？}
与原始 Transformer（Encoder-Decoder）和 BERT（Encoder-only）不同，GPT 舍弃了 Encoder，只保留 Decoder。
\begin{enumerate}
    \item \textbf{生成任务的天性}：生成任务是自左向右的，Decoder 的 Masked Self-Attention 完美契合这种训练方式。
    \item \textbf{自回归建模}：GPT 将条件概率分解为 $P(x) = \prod P(x_i | x_{<i})$。这种因果建模使得模型在推理时能够逐词生成。
\end{enumerate}

\subsection{结构组件}
一个典型的 GPT Block 包含：
\begin{itemize}
    \item \textbf{Masked Multi-Head Attention}：确保当前计算只能看到之前的词，防止“偷看答案”。
    \item \textbf{Layer Normalization}：通常采用 Pre-LN 结构以增加深层模型的稳定性。
    \item \textbf{Position-wise FFN}：两层全连接网络，通常使用 GeLU 激活函数。
\end{itemize}

\section{Evolution: 从 GPT-1 到 GPT-4}

\begin{table}[h]
\centering
\begin{tabularx}{\textwidth}{|l|l|X|}
\hline
\textbf{模型} & \textbf{参数量} & \textbf{核心贡献} \\
\hline
GPT-1 & 1.17 亿 & 确立了“生成式预训练 + 判别式微调”的二阶段流程。 \\
\hline
GPT-2 & 15 亿 & 强调 \textbf{Zero-shot} 能力，展示了模型在不经微调下处理多任务的潜力。 \\
\hline
GPT-3 & 1750 亿 & 引入 \textbf{In-Context Learning (ICL)}，通过 Few-shot 提示即可完成复杂任务。 \\
\hline
InstructGPT & 1750 亿 & 引入 \textbf{RLHF}，使模型输出符合人类指令和偏好（Alignment）。 \\
\hline
GPT-4 & 未公开 (万亿级) & 多模态理解力、处理长输入、极强的推理和逻辑能力。 \\
\hline
\end{tabularx}
\end{table}

\section{RLHF：对齐人类偏好}
为了让模型不产生胡言乱语或有害内容，OpenAI 引入了\textbf{人类反馈强化学习（RLHF）}：
\begin{enumerate}
    \item \textbf{监督微调 (SFT)}：在人工编写的问答对上进行微调。
    \item \textbf{奖励模型 (RM) 训练}：人类对模型的多个输出进行排序，训练一个评分模型。
    \item \textbf{强化学习 (PPO)}：使用奖励模型作为反馈，通过强化学习优化策略，使得生成内容得分最高。
\end{enumerate}

\section{GPT vs. BERT}
\begin{table}[h]
\centering
\begin{tabularx}{\textwidth}{|l|X|X|}
\hline
\textbf{特性} & \textbf{BERT} & \textbf{GPT} \\
\hline
\textbf{架构} & Encoder-only & Decoder-only \\
\hline
\textbf{注意力} & Bidirectional (双向) & Unidirectional (因果/单向) \\
\hline
\textbf{预训练目标} & MLM (填空) & CLM (预测下一个词) \\
\hline
\textbf{擅长领域} & 理解、分类、抽取 & 创意写作、对话、代码生成 \\
\hline
\end{tabularx}
\end{table}

\section{关键代码：Causal Self-Attention mask}
GPT 的核心在于其因果掩码，在 PyTorch 中通常这样实现：
\begin{lstlisting}
import torch
import torch.nn as nn

def create_causal_mask(size):
    # 生成一个上三角矩阵，并取反，使得左下角(包含对角线)为 1
    mask = torch.tril(torch.ones(size, size))
    return mask == 0 # True 的地方将被填充为 -inf

# 在 Attention 计算中的应用
# scores: [Batch, Heads, Seq, Seq]
mask = create_causal_mask(seq_len)
scores = scores.masked_fill(mask, float('-inf'))
attn = torch.softmax(scores, dim=-1)
\end{lstlisting}

\section{深度面试题}

\textbf{问：为什么 GPT 相比 BERT 展现出了更强的通用泛化能力？}\\
答：1. \textbf{预训练难度}：CLM（预测未来）比 MLM（填补中间）更难，模型被迫学习更深层的逻辑和语义；2. \textbf{任务统一}：自然语言的所有任务都可以转化为生成任务（即 Prompting），这使得 GPT 可以直接进行 Zero-shot 迁移，而 BERT 往往需要针对每个下游任务修改 Head 并全量微调。

\textbf{问：什么是涌现能力 (Emergent Abilities)？}\\
答：指当模型参数量突破某个临界点（如 600 亿）时，模型在某些任务（如逻辑推理、数学计算）上的表现会突然从接近随机猜测跃升到远超标准水平。这是大模型（LLM）区别于小模型的最显著特征。

\section{参考文献}
\begin{enumerate}
    \item Radford, A., et al. (2018). \textit{Improving Language Understanding by Generative Pre-Training} (GPT-1).
    \item Brown, T., et al. (2020). \textit{Language Models are Few-Shot Learners} (GPT-3).
    \item Ouyang, L., et al. (2022). \textit{Training language models to follow instructions with human feedback} (InstructGPT).
\end{enumerate}

\end{document}
